\subsubsection{Counting the classes of the reference chain}
\label{subsec:selected-class-counts}
\hypertarget{source-selected-class-counts}{}

The reference norm is a sum over definite classes. Bounding the number of its terms requires a class count for the ranks occurring along this chain.

Fix the reference tuple, the prime assigned to the sequence \(T=T_u\), and the two
type primes \(g_0,g_1\) for the chosen CM extension \(E/F\). Write
\(O=\mathcal O_E\), \(d=[E:\mathbb Q]=2f\), and
\(r=|D_E|^{1/d}\). Let \(n\) be a rank occurring in the count,
\(1\le n\le N_m\). We first bound the number of rational unitary
orbits of self-dual \(O\)-lattices in the standard positive Hermitian
space \((E^n,h)\).
\label{cor:selected-class-counts-1}

\begin{proposition}[Size of the reference class data]\label{prop:reference-class-data}
For the selected reference tuple on \(E/F\), with type primes
\(g_0,g_1\) and assigned prime \(T_u\),
use the logarithmic size unit \(U_u=fm^2N_m^2\log(g_0g_1T_u)\).
Write \(h,i,s\) for the class, matrix-index and intermediate-coset counts. Then
\[
 h\le e^{55U_u},\qquad i\le e^{U_u},\qquad s\le e^{2U_u}.
\]
 With the
metrics of the selected chain, their logarithmic squared-norm
costs are at most \(112U_u\) in an ordinary step and
\(225U_u\) in the two coupled end steps together.
\end{proposition}

\begin{proof}
\proofstep{Find a short basis by transference against the trace dual.}
Give a lattice \(M\) the real inner product
\(\operatorname{Tr}_{E/\mathbb Q}h(x,y)\), with the convention that
the trace sums over all embeddings. Its integral trace dual is
\(\mathfrak D_E^{-1}M\). Indeed the trace pairing is nondegenerate
on the rational space. Testing a candidate dual vector against all
\(O\)-multiples of every vector of \(M\) puts its Hermitian
pairings in the codifferent, and localization of Hermitian
self-duality gives the displayed module.

If \(x\ne0\) belongs to
this dual, \(h(x,x)\) is totally positive and belongs to
\(\mathfrak D_E^{-2}\), since the different is conjugation stable.
Thus \(|N_{E/\mathbb Q}h(x,x)|\ge |D_E|^{-2}\). Since \(d=[E:\Q]\) and \(r=|D_E|^{1/d}\), applying the arithmetic
and geometric mean inequality to the conjugates gives
\[
 \|x\|^2=\operatorname{Tr}_{E/\Q}h(x,x)
     \ge d\,|N_{E/\Q}h(x,x)|^{1/d}\ge d/r^2.
\]
The Euclidean transference theorem in dimension \(dn\) bounds the
last successive minimum of \(M\) by \(n\sqrt d\,r\). These minima
are attained by discreteness. Choose an \(E\)-basis \(v_1,\ldots,v_n\)
from a rationally spanning family attaining this bound.
\cite{banaszczyk}

\proofstep{Count Gram matrices and the lattices between two integral bounds.}
Put \(B=4f^3 n^4 r^2\). The Gram matrix \(A=(h(v_i,v_j))\) has
entries in \(O\). Cauchy--Schwarz at each embedding bounds their
absolute values by \(dn^2r^2\le B\). The unit ball for the maximum
of the archimedean absolute values has real dimension \(d\).
Distinct elements of \(O\) have distance at least one in this norm
by the product formula. Packing translates of its open half ball
bounds the number of choices per entry by \((1+2B)^d\).
In the \(v_i\) basis, with the form written \(\bar x^t A y\),
\[
           O^n\subseteq M\subseteq A^{-1}O^n .
\]
The order of the quotient of the outer modules is
\(|N_{E/\mathbb Q}\det A|\le(nB)^{dn}\) by Hadamard.

Each submodule of this quotient is generated by at most \(n\)
elements. Locally its inverse image in \(A^{-1}O_P^n\) is a
lattice free of rank \(n\) over the DVR. Pad the lists of
generators of its images to length \(n\), and combine them at the
finitely many primes of the quotient using the Chinese remainder
theorem. If \(H=A^{-1}O^n/O^n\), the generator lists give a surjection
\[
 H^n\longrightarrow\{U:O^n\subseteq U\subseteq A^{-1}O^n\},
 \qquad (x_1,\ldots,x_n)\longmapsto
  \text{the inverse image of }\sum_i O x_i .
\]
Hence the count is at most \(|H|^n\le(nB)^{dn^2}\).
The pair consisting of \(A\) and this coordinate module determines
one rational unitary class, by mapping an ordered basis to another
basis with the same Gram matrix.

Set \(P=d^2 n^4 r^2\), so \(B=fP\). As \(B=fP\ge1\) and \(f^2n\le P\), the matrix count times the
module count is bounded by
\[
 (1+2B)^{dn^2}(nB)^{dn^2}
    \le (3f^2nP^2)^{dn^2}\le(3P)^{3dn^2}.
\] This
argument allows projective \(O\)-lattices of every Steinitz class.

At a selected place \(v\) over \(T\) a decoration specifies a
residue functional. At a selected place over a type prime it
specifies \(n\) such functionals, the coordinates of its frame.
On a fixed lattice of rank \(n\), each functional is determined by
its values on lifts of a residue basis. The localized lattice
modulo its maximal ideal has dimension \(n\). One can take the
lifts in the global lattice since denominators outside the place
are invertible on the residue field. The rules for a decoration
determine it by restriction to the localized lattices, including
when its definition represents the functions on the ambient space.

Writing \(q_v\) for the residue cardinality, we have at most
\[
        \prod_{v\in S_T}q_v^n\prod_{v\in S_G}q_v^{n^2}
          \le (g_0g_1T)^{fn^2}
\]
choices. The selected half above each of the three split primes has
\(f\) places of residue degree one. Here \(S_T\) can be a subset
of the word half and \(S_G\) is the union of the two type halves.
Transport each decorated point to one of the underlying lattice
representatives by a rational unitary transformation, leaving the
external frame-change group element equal to the identity. For a fixed underlying lattice each class costs at most the decoration count just obtained. Summing gives
\[
 \#\{\text{principal decorated classes}\}\le \#\{\text{underlying lattice classes}\}\,\max_L\#\{\text{decorations on }L\}.
\] Stabilizers can only reduce this number.

We explain how the same bound applies to each of the three class
spaces in the period chain. Suppose a point space injects into the
decorated points and every action of \((g,1)\), for \(g\) rational
unitary, on an image point lifts to an allowed move. Choose one
representative per orbit in the point space. If two chosen images
belong to one principal class, lift the relating move and apply
injectivity. Concretely, for the injection \(i\) on points and its allowed lifted
move \(\tilde g\), one has
\[
 i(y)=(g,1)i(x)=i(\tilde g x)\quad\Longrightarrow\quad y=\tilde g x .
\]
Thus taking principal classes does not identify distinct point-space
orbits.

In a class space of
a column of rank \(N_j\), forgetting the supplementary twist
labels is injective because the base tuple has empty twist
support. The move lifts as \((g,1,1)\) at every matrix index.
For the evaluation rows take rank \(N_m-1\), empty word decoration,
and the identity on decorated points. The move is \((g,1)\).
For the top class count take rank \(N_m\), the full word decoration
and the identity, with \((g,1)\) in the permitted product subgroup.
The two type halves exhaust the row places at the base tuple. Thus
all three counts use precisely the prime product \(G=g_0g_1T\).

The selected size conditions give the following simultaneous range:
\[
 L\ge1,\qquad\log r\le L,\qquad fN_m\le e^L\le g_0 .
\] In this comparison it is
enough to use these four inequalities. Since \(12\le e^{3L}\),
\[
         3P=12(fn)^2 n^2 r^2\le e^{9L}\le G^9 .
\]
The resulting bound on each of the three class counts is
\[
 G^{(54+1)fn^2}\le\exp(55U_u),\qquad U_u=fm^2N_m^2\log G .
\] This is uniform in the base row,
column, index and word. \label{cor:selected-class-counts-2}

\proofstep{Read each index on the residue bases of the padded lattices.}
For the matrix indices only stages \(j+2\le m\) occur. The
initial index is unique. For a positive stage record the residual
matrix of its representative at each type place on the padded
lattice of rank \(n'=N_j+c_j\le N_m\). In coordinates with lattice
exponents \(e_a\) (zero on the initial block and successive
positive integers on the padding), its entry is the reduction of
\(g_i^{e_b-e_a}h_{ab}\) at a place over \(g_i\).

In these coordinates the residual matrix is the reduction of
\(D^{-1}hD\), with \(D=\operatorname{diag}(g_i^{e_a})\). Equal
residual matrices for \(h_1,h_2\) imply
\[
 D^{-1}h_1^{-1}h_2D
     =(D^{-1}h_1D)^{-1}(D^{-1}h_2D)\equiv I\pmod{g_i}.
\] The type prime is unramified of residue degree one.
Every defining flag lattice lies between the padded lattice and
its \(g_i\)-multiple, and the prescribed residue vectors lie in the
padded lattice. Hence the relative action preserves each flag
lattice and fixes the vector classes in its prescribed quotients.
This is the right stabilizer in the local double coset. Equality
of these double cosets at all the selected places identifies the
index by its defining injectivity. At the base tuple the resulting count of indices satisfies
\[
 \#\{\text{indices}\}\le(g_0g_1)^{f(n')^2}\le e^{U_u}.
\]

\proofstep{Bound the cosets by the single modulus at the two type primes.}
The intermediate coupling set consists of cosets of a module
\(A_0\) represented by elements of a module \(A_1\subset E^n\).
The defining local conditions give \(g_0g_1 A_1\subseteq A_0\):
at a place over \(g_i\) the one relaxed lattice exponent is
restored by multiplication by \(g_i\), and there are no twist
places. To count the quotient, let \(L\subset A_1\) be generated by
finitely many representatives. The inclusion just proved gives
\[
 L/(g_0g_1)L\twoheadrightarrow L/(L\cap A_0),\qquad
 |L/(L\cap A_0)|\le (g_0g_1)^{\operatorname{rank}_{\mathbb Z}L}
                 \le (g_0g_1)^{2fn}.
\]
Every finite set of cosets is counted in such an image. This proves the bound even for the semilocal modules.

The coset bound applies to every actual intermediate lattice as well as on the
standard one. We may take \(C_{\rm cos}=2\),
\(C_{\rm idx}=1\), and \(C_{\rm cl}=55\). In particular the
ordinary and coupled squared-norm losses above are at most
\(112 U_u\) and \(225 U_u\), respectively. Increasing the
single-kernel constant by \(58\) pays \(116U_u\) per lift.
Both end lifts, at the same word and embedding, are charged
in the coupled case. \label{cor:selected-class-counts-3}
\end{proof}

